% ============================================================
% Chapter III — Methodology
% ============================================================
\section{Methodology}

This chapter outlines the six-stage methodological framework used to
develop, validate, and interpret the multivariate predictive model. The
framework follows Hair et al.'s\parencite{hair2019} guidelines for
dependence multivariate techniques, adapted for the specific characteristics
of Formula One telemetry data.

\subsection{Stage 1: Define the Research Problem and Technique Selection}

The objective is to predict a single multi-class non-metric dependent variable
--- \textbf{Finishing Tier} (Podium, Point Scorer, No Points) --- using
multiple metric and non-metric independent variables derived from lap-level
telemetry, race conditions, and strategy decisions.

Two dependence techniques were considered:

\begin{description}
  \item[Multiple Discriminant Analysis (MDA)] The theoretical baseline for
    multi-group classification. MDA would estimate canonical discriminant
    functions to maximize between-group variance. However, MDA requires
    homogeneity of covariance matrices and multivariate normality ---
    assumptions that are routinely violated in motorsport data due to the
    skewed nature of track evolution and safety car deployments.

  \item[Multinomial Logistic Regression] The definitive technique selected
    for this study. Using Maximum Likelihood Estimation (MLE), it calculates
    the specific log-odds of a driver landing in each tier relative to a
    baseline category. It offers greater operational value for risk
    assessment than MDA's rigid classification boundaries and makes no
    assumption about covariance structure.
\end{description}

MDA serves as the benchmark against which the Multinomial Logistic Regression
is compared, but the latter is the primary inferential engine.

\subsection{Stage 2: Develop the Analysis Plan (Research Design)}

The research design encompasses the analytical decisions made before model
estimation.

\subsubsection{Categorical Encoding}

Qualitative race strategy variables (tire compounds: Soft, Medium, Hard,
Intermediate, Wet) are non-metric and cannot enter a regression equation
directly. They were converted into $k{-}1$ dummy variables, where $k$ is the
number of compound categories. The first alphabetical compound (Hard) was
dropped as the baseline reference, yielding binary indicators for
\texttt{Tire\_SOFT} and \texttt{Tire\_MEDIUM}. This avoids the dummy variable
trap while assigning specific mathematical weights to each compound relative
to the reference.

\subsubsection{Sample Size}

The processed dataset contains 2,880 lap-level observations across 6
independent variables. This yields 480 observations per predictor --- well
above the recommended minimum of 15 to 20 observations per independent
variable for reliable multivariate estimation\parencite{hair2019}.

\subsubsection{Structural Break Classification (The DNF Protocol)}

Drivers who retired mid-race (crashes, mechanical failures) represent a
\textit{structural break} --- not a standard missing data point to be
imputed or dropped. A DNF lap carries telemetry from a driver who may have
been running P2 pace before a suspension failure, but whose final
classification is ``Retired.'' Including these laps in the Podium cluster
would contaminate the tier with telemetry from a driver who never finished.

The protocol was:
\begin{enumerate}
  \item Flag all drivers with \texttt{Status = "Retired"} or
    \texttt{ClassifiedPosition = "R"}.
  \item Reassign \textit{all} their completed laps to the ``No Points'' tier.
  \item This preserves the integrity of lap-time telemetry by ensuring that
    only drivers who actually finished in a given tier contribute to that
    tier's covariance structure.
\end{enumerate}

\subsection{Stage 3: Evaluate the Assumptions Underlying the Techniques}

\subsubsection{The Methodological Pivot: Covariance Diagnostics}

Box's M test was conducted to assess the equality of covariance matrices
across the three finishing tiers. The test yielded
$\chi^{2}(20) = 160.419$, $p < .001$, formally rejecting the null
hypothesis of equal covariance structures.

This empirical failure --- driven by the heavily skewed distribution of
Podium laps (14.9\%) versus No Points laps (51.0\%), combined with the
heterogeneous variance introduced by track evolution and safety car
deployments --- formally justifies abandoning MDA as the primary technique.
Multinomial Logistic Regression was selected specifically because it
bypasses this rigid normality requirement.

\subsubsection{Coupled Variables and Dynamic Interaction Terms}

Tire degradation and fuel mass burn-off are physically coupled variables:
both increase monotonically with lap number. Initial VIF diagnostics
confirmed multicollinearity risk between these variables and their product
term. Following Hair et al.'s\parencite{hair2019} recommendation, the
variables were \textit{not} dropped. Instead, a dynamic interaction term
was constructed:

\[
  \text{Degradation} \times \text{Fuel} =
  (\text{Tire Age} - \bar{x}_{\text{Tire Age}}) \times
  (\text{Fuel Mass} - \bar{x}_{\text{Fuel Mass}})
\]

Mean-centering before multiplication eliminates structural
multicollinearity while preserving the interaction's ability to capture
how these forces compound each other over the course of a stint. Post-
interaction VIF scores confirmed that all predictors remained below the
threshold of 10.

\subsection{Stage 4: Estimate the Multivariate Model and Assess Overall Fit}

\subsubsection{Model Estimation}

The Multinomial Logistic Regression was estimated via Maximum Likelihood
Estimation (MLE) using the \texttt{multinom()} function in R's
\texttt{nnet} package. The model computes simultaneous log-odds equations
for the Podium and Point Scorer tiers relative to the No Points baseline:

\[
  \ln\!\left(\frac{P(Y = k)}{P(Y = \text{No Points})}\right) =
  \beta_{0k} + \beta_{1k} \cdot \text{Grid} +
  \beta_{2k} \cdot \text{TireAge} +
  \beta_{3k} \cdot \text{Fuel} +
  \beta_{4k} \cdot \text{Deg}\times\text{Fuel} +
  \beta_{5k} \cdot \text{Tire\_MED} +
  \beta_{6k} \cdot \text{Tire\_SOFT}
\]

where $k \in \{\text{Podium}, \text{Point Scorer}\}$.

\subsubsection{Model Fit Assessment}

Overall fit was evaluated using pseudo-$R^2$ equivalents (McFadden's
$R^2$) derived from the residual deviance and AIC. The classification hit
ratio was computed to confirm that the model places drivers into their
actual finishing tiers at a rate significantly higher than proportional
chance (33.3\% for a three-class problem with equal priors, or 51.0\% for
a majority-class baseline).

\subsection{Stage 5: Interpret the Variate(s)}

\subsubsection{Identify Key Drivers}

The MLE coefficients (log-odds) were converted to odds ratios
($e^{\beta}$) to isolate which telemetry metrics and strategies exert the
strongest influence on podium versus no-points outcomes. The sign and
magnitude of each coefficient indicate whether a predictor increases or
decreases the probability of a given tier.

\subsubsection{VIF Validation Rule}

Before interpreting any coefficient, VIF scores were verified to be below
10 to ensure that standard error inflation was controlled. This prevents
sign reversals --- a well-documented artifact of multicollinearity where a
beneficial race strategy could mathematically appear harmful in the
coefficient estimate\parencite{tabachnick2019}. Coefficients were compared
on their raw (unstandardized) scale since all metric predictors are
already on comparable physical units (laps, kilograms).

\subsection{Stage 6: Validate the Multivariate Model}

\subsubsection{Generalizability via Leave-One-Group-Out Cross-Validation}

To prove the model's predictive power outside the training sample,
Leave-One-Group-Out Cross-Validation (LOGO-CV) was deployed using
\texttt{Race\_ID} as the grouping parameter. In each iteration:

\begin{enumerate}
  \item The model was trained on $N{-}1$ Grand Prix events.
  \item The held-out race's laps were predicted without ever having been
    seen during training.
  \item Predictions were aggregated across all three folds to produce an
    overall classification report (precision, recall, F1-score, accuracy).
\end{enumerate}

To prevent intra-race data leakage, the interaction term was mean-centered
within each training fold (not on the full dataset) before application to
the test fold. This ensures the model cannot exploit statistical artifacts
that only exist when all races are combined.

A limitation of this validation is that non-randomized structural changes
--- such as year-to-year FIA regulation shifts, tire compound formula
changes, or major aerodynamic overhauls --- are not captured. The model
generalizes across \textit{circuit architecture} but not across
\textit{regulatory eras}.

\subsubsection{Implementation Architecture}

The analytical pipeline was executed across two environments:

\begin{itemize}
  \item \textbf{Python (pandas, scikit-learn):} Data wrangling, dummy
    encoding, interaction term generation, and LOGO-CV validation loops.
    The \texttt{LogisticRegression} class with \texttt{class\_weight=
    "balanced"} was used for cross-validation.
  \item \textbf{R (\texttt{nnet}, \texttt{car}, \texttt{heplots}):}
    Theoretical proof via \texttt{multinom()} for rigorous coefficient
    summaries, exact $p$-values, Box's M test, and VIF diagnostics.
\end{itemize}

\subsection{Sampling Strategy}

The unit of analysis is the individual lap. Every completed lap from every
driver across the 2023 Bahrain, Saudi Arabian, and Australian Grands Prix
was included, yielding 2,880 observations after DNF reassignment and
missing-value removal. No random sampling was performed; the dataset
represents the complete population of laps from the three selected races.
Races were chosen purposively to cover three architecturally distinct
circuit types: permanent (Bahrain), street (Jeddah), and parkland (Albert
Park).

\subsection{Tools and Instruments}

\begin{table}[htbp]
  \centering
  \caption{Software and Libraries}
  \label{tab:tools}
  \begin{tabular}{lll}
    \toprule
    \textbf{Tool} & \textbf{Version} & \textbf{Purpose} \\
    \midrule
    Python & 3.11 & Data preparation and ML validation \\
    FastF1 & 3.8.2 & F1 telemetry data extraction \\
    pandas & 2.3.3 & Data manipulation and encoding \\
    scikit-learn & 1.8.0 & Logistic regression, LOGO-CV \\
    R & 4.x & Statistical inference \\
    nnet (R) & --- & \texttt{multinom()} for MLE estimation \\
    car (R) & --- & VIF diagnostics \\
    heplots (R) & --- & Box's M test \\
    \bottomrule
  \end{tabular}
\end{table}

\subsection{Data Collection Process}

Data were obtained from the FastF1 Python API, which interfaces with the
official Formula One timing and telemetry data feed. For each race:

\begin{enumerate}
  \item The race session was loaded via \texttt{fastf1.get\_session()} with
    telemetry, lap, and weather data enabled.
  \item All laps were concatenated into a single dataframe and tagged with a
    \texttt{Race\_ID} for group-aware cross-validation.
  \item The \texttt{Finishing\_Tier} target was mapped from final race
    classification positions.
  \item DNF drivers were flagged and their laps reassigned to ``No Points.''
  \item Engineered features (fuel mass proxy, tire age, interaction term)
    were computed and dummy-encoded.
\end{enumerate}

The final dataset was exported as \texttt{f1\_clean\_data.csv} with 2,880
rows and 10 variables.

\subsection{Ethics and Compliance}

This study uses publicly available telemetry data from the FastF1 API, which
interfaces with the official Formula One data feed. No human participants
were involved, and no personally identifiable information beyond public
driver abbreviations (e.g., ``VER'' for Max Verstappen) was collected. The
study complies with FastF1's terms of use for academic research.

\subsection{Ensuring Credibility}

\begin{itemize}
  \item \textbf{Data transparency:} The full dataset and all source code
    (Jupyter notebooks, R scripts, figure generation) are available in the
    project repository for independent verification.
  \item \textbf{Assumption testing before estimation:} Box's M test and VIF
    diagnostics were conducted and reported regardless of outcome, following
    the sequential framework.
  \item \textbf{Strict validation design:} LOGO-CV was chosen over random
    train-test splits to prevent data leakage and produce honest out-of-sample
    performance estimates. Interaction terms were re-centered within each
    training fold.
  \item \textbf{Balanced evaluation metrics:} Class weighting was applied
    during validation, and performance is reported via precision, recall,
    and F1-score --- not just accuracy --- to prevent misleading claims on
    imbalanced data.
  \item \textbf{Dual-environment cross-verification:} Statistical inference
    (R) and predictive validation (Python) were executed independently, and
    results were cross-verified to ensure consistency across implementations.
\end{itemize}
